% Appendix, placed after the references by main.tex. The course lists the
% contribution statement as a deliverable separate from the 5--8 page report
% and requires the LLM usage statement to be in the report; confirm with the
% teaching team that an appendix does not count toward the page limit.
\section{Contribution Statement}
% Drafted 2026-10-12 from two sources: the group's task-assignment and
% presentation sheet (who was assigned which task and track, who presents which
% part), and what is visible in the repository (reports, commits, pull
% requests). Rows whose evidence is TBD rest on the assignment sheet alone.
% Chu Jiajie's row was given by the group on 2026-10-13: the PegInsertionSide
% backbone runs follow the assignment sheet (the task itself is left out of the
% report), and the PlaceSphere demonstration data is credited to him; the
% repository shows no commits or report under his name for either.
% One split follows the assignment sheet, as decided by the group on
% 2026-10-12: the PickCube split between Wong Pak Hung Philip (data size) and Yang Haoran
% (backbone); the PickCube report and its run directory, covering both tracks,
% are under Yang Haoran's account.
% Every member should confirm or correct their own row.
\todo{Each member confirms their row.}
{\small
\begin{tabularx}{\linewidth}{@{}p{0.17\linewidth}YY@{}}
  \toprule
  Member & Contribution & Supporting evidence \\
  \midrule
  Chu Jiajie & Backbone comparison (Transformer and MLP) on PegInsertionSide,
  a task not included in this report; generation of the PlaceSphere
  demonstration data and the PlaceSphere data-size study. & Group assignment
  sheet; PlaceSphere data and results recorded in the repository. \\
  Jin Qizheng & StackCube data-size study ($N=25$ to 400) and its report;
  provided the $N=100$ UNet baseline runs for the StackCube backbone
  comparison. & \texttt{report\_stackcube\_scaling.md},
  \texttt{handoff\_lusen\_n100.md}; pull requests 9 and 14 of the integration
  repository. \\
  Lu Sen & StackCube backbone comparison; debugging of the experiment
  environment and setup of the run framework on the Moore Threads server;
  early development of the training and evaluation framework; presents the
  data-size and backbone studies. &
  Member report \texttt{task2\_backbone}; launch script
  \texttt{moer\_train.sh} for the Moore Threads server; presentation plan. \\
  Shao Chuncheng & LiftPegUpright data-size study. &
  Member report \texttt{task6}. \\
  Song Zirui & LiftPegUpright backbone comparison; cluster job template; early
  development of the training and evaluation framework; presents demonstration
  generation, policy setup, and baselines. & Member report
  \texttt{task6\_backbone}; pull request 1 of the integration repository;
  presentation plan. \\
  Wang Bowen & PushCube data-size study and backbone comparison. & Member
  report \texttt{task3}. \\
  Wong Pak Hung Philip & PickCube data-size study; presents the research
  questions and task selection. & Group assignment sheet; presentation plan. \\
  Yang Haoran & PickCube backbone comparison with cost measurements, and the
  PickCube report; Moore Threads reproduction and the attention fix; early
  framework development; maintenance of the integration repository. &
  \texttt{report\_pickcube\_yhr.md}; branch
  \texttt{moore-transformer-attn-fix}; integration pull requests. \\
  Yang Qihang & PullCube data-size study and backbone comparison (26 runs). &
  Member report \texttt{report\_pullcube\_yqh.md}. \\
  Yu Houcheng & Training and evaluation pipeline (\texttt{dp-manip}) and
  demonstration generator; failure-guidance and takeover-correction
  studies; cross-task results summary; presents the failure-correction study.
  & Commit history of \texttt{dp-manip} and \texttt{maniskill-demogen};
  experiment notes under \texttt{dp-manip/docs}; presentation plan. \\
  \bottomrule
\end{tabularx}}

\section{LLM Usage Statement}
% The course policy asks for the tools used, the sections or components
% affected, and how the tools contributed. A fourth column for records or
% verification was dropped on 2026-10-13 because most rows had nothing to cite.
{\small
\begin{tabularx}{\linewidth}{@{}p{0.26\linewidth}p{0.2\linewidth}Y@{}}
  \toprule
  Section / component & Tool & Contribution \\
  \midrule
  Report scaffold & OpenAI Codex; Claude Code & Section structure, placeholder
  tables, cross-check against the course requirements, compilation and layout
  checks \\
  Abstract, Sections 1--6, Tables 1--5, Figure 2 & Claude Code & First draft
  of the text from the repository's configurations, code, and experiment
  notes; transcription of table values from the group's results summary and
  member reports; plotting script; bootstrap intervals for PlaceSphere in
  Section 4.1; lookup of reference metadata \\
  Training and evaluation code & Claude & Generation of the cluster job
  scripts; reorganization of the code; optimization of the data loader for
  cluster training \\
  Single-machine framework & WorkBuddy & Generation of a single-machine
  training framework used to validate the pipeline \\
  Moore Threads environment (Section 4.4) & OpenAI Codex & Assistance with
  configuring the software environment on the Moore Threads server \\
  Experiment data on the servers & DeepSeek & Transfer and organization of
  the data on the remote servers \\
  Presentation videos & Claude & Generation of the videos shown in the
  presentation \\
  Presentation slides & Claude (Artifacts) & Assistance with preparing the
  slides for the final presentation \\
  \bottomrule
\end{tabularx}}

\smallskip\noindent Commits made with Claude carry a co-author trailer in the
repository history.

\section{Reproducibility and Acknowledgments}
% Repository names are from the top-level README of the integration repository.
% All three were public on GitHub when checked on 2026-10-12.
Code is public in three repositories. Every run records its resolved
configuration, selected demonstration seeds, and normalization statistics.
\begin{itemize}
  \item Course integration repository, including this report:
  \url{https://github.com/hryang1130/7606C}
  \item Training and evaluation pipeline:
  \url{https://github.com/hollinsStuart/dp-manip}
  \item Demonstration generator:
  \url{https://github.com/hollinsStuart/maniskill-demogen}
\end{itemize}
\todo{Give the commit of each repository used for the reported runs, and links
to the generated datasets, the trained checkpoints, and a six-task video index
covering expert generation and policy rollouts. Verify that the links are
accessible to the teaching team.}

We use ManiSkill and its motion-planning solvers~\citep{tao2024maniskill3}.
The UNet denoiser is the implementation shipped with the ManiSkill Diffusion
Policy baseline, which derives from the Diffusion Policy
code~\citep{chi2023diffusion}; the Transformer denoiser follows the
transformer variant of the same code; and the noise scheduler comes from the
\texttt{diffusers} library. Experiments ran on the HKU CS GPU Farm and on a
Moore Threads server provided for the course.
